\section{Reconstruction implies capacity}

The necessity argument has a useful form that does not require a uniform
contraction bound. For arbitrary parameters $0<\lambda_e\le1$, retain
\eqref{eq:rho}, put
\[
 \omega_e=1-\lambda_e^4,
\]
and define the edge-noise-weighted capacity
\begin{equation}\label{eq:weighted-capacity}
 \Cap_3^{\mathrm{noise}}(T,\rho)
 =\sup\left\{|\theta|:
 \sup_{\xi\in\partial T}
 \sum_{v\in\xi\setminus\{o\}}
 \omega_{e_v}\left(\frac{\theta(e_v)}{\rho(v)}\right)^2\le1
 \right\}.
\end{equation}

\subsection{A cubic statistic removes degree dependence}

Set
\begin{equation}\label{eq:f}
 K=\frac c3=\frac1{3000},\qquad f(t)=t+Kt^3,
 \qquad D_f=1+9K.
\end{equation}
If $t=m(\mu\star\nu)$ for laws in $\calC$, then \eqref{eq:input} and
$t\le u+v$ give
\begin{align*}
 f(t)&\le u+v-cuv(u+v)+K(u+v)^3\\
 &=f(u)+f(v)+(3K-c)uv(u+v)=f(u)+f(v).
\end{align*}
Thus $f(m)$ is subadditive under any finite product. The number of children
never enters the constant.

Let $m_v^{(n)}$ be the information at $v$ from descendants at distance $n$,
and set $m_v=\lim_n m_v^{(n)}$. Data processing gives these limits in
$[0,3]$. For $r_{vw}=\lambda_{vw}^2$, channel scaling, subadditivity of
$f(m)$, and passage to the limit imply
\begin{equation}\label{eq:fy}
 f(m_v)\le\sum_{w\in\ch(v)}f(r_{vw}m_w).
\end{equation}
All sums are finite. No closure assertion about a limiting posterior is used.

Put $y_v=f(m_v)$, set $\omega_{vw}=1-\lambda_{vw}^4$, and define
\[
 \delta_{vw}=\frac{K(1-r_{vw}^2)}{D_f^3}
 =\frac{K\omega_{vw}}{D_f^3}.
\]
Since $y_w\le D_fm_w$ and
\[
 f(rm)=rf(m)-Kr(1-r^2)m^3,
\]
equation \eqref{eq:fy} yields
\begin{equation}\label{eq:edge-concave}
 y_v\le\sum_w r_{vw}y_w(1-\delta_{vw}y_w^2)
 \le\sum_w\frac{r_{vw}y_w}
                    {\sqrt{1+2\delta_{vw}y_w^2}}.
\end{equation}
The last inequality uses $1-t\le(1+2t)^{-1/2}$ for $t\ge0$.

\subsection{The edge-weighted flow}

Choose coefficients
\[
 0\le d_{vw}\le
 \frac{r_{vw}}{\sqrt{1+2\delta_{vw}y_w^2}},
 \qquad y_v=\sum_w d_{vw}y_w.
\]
They are obtained by multiplying all upper-bound coefficients by the ratio
of $y_v$ to the right side of \eqref{eq:edge-concave}; if that side is zero,
use zero coefficients. Define
\[
 \ell_o=1,
 \qquad \ell_w=\ell_vd_{vw},
 \qquad \vartheta(v,w)=\ell_wy_w.
\]
These equations give flow conservation and root mass $y_o=f(m_\infty)$.
With $g_v=\ell_v/\rho(v)$ and $\rho(w)=\rho(v)r_{vw}$,
\[
 g_w^2(1+2\delta_{vw}y_w^2)\le g_v^2,
 \qquad
 \frac{\vartheta(v,w)}{\rho(w)}=g_wy_w.
\]
Therefore
\[
 \omega_{vw}\left(\frac{\vartheta(v,w)}{\rho(w)}\right)^2
 \le \frac{D_f^3}{2K}(g_v^2-g_w^2).
\]
Summing along a ray telescopes. After scaling the flow, we obtain the
following quantitative obstruction.

\begin{theorem}[Weighted-capacity necessity]\label{thm:weighted-necessity}
For an arbitrary locally finite deterministic rooted tree with
$0<\lambda_e\le1$, reconstruction implies
\begin{equation}\label{eq:weighted-necessity}
 \Cap_3^{\mathrm{noise}}(T,\rho)
 \ge \sqrt{\frac{2K}{D_f^3}}\,f(m_\infty)
 \ge \sqrt{\frac{2K}{D_f^3}}\,m_\infty>0.
\end{equation}
\end{theorem}

Under the uniform bound \eqref{eq:channels}, one has
$\omega_e\ge1-q^4$. Hence a weighted-admissible flow scaled by
$\sqrt{1-q^4}$ is admissible for \eqref{eq:capacity}. Consequently,
\begin{equation}\label{eq:necessityquant}
 \Cap_3(T,\rho)
 \ge\sqrt{2\delta_q}\,f(m_\infty)
 \ge\sqrt{2\delta_q}\,m_\infty.
\end{equation}
This proves the necessity implication in Theorem~\ref{thm:main}. It also gives
$m_\infty\le\Cap_3(T,\rho)/\sqrt{2\delta_q}$. Since $m_\infty\le3$ and
$(2\delta_q)^{-1/2}>3$, this is the upper bound in
\eqref{eq:quantitative-m}. The upper total-variation bound in
\eqref{eq:quantitative} follows from $a_\infty\le\tfrac12\sqrt{m_\infty}$.
